\section{Introduction}
\label{sec:introduction}

In robot racing, the course is often presented as background: a map on which a car is
controlled, a sequence of gates through which a quadrotor flies, or a set of buoys
around which a vessel navigates. For learning and evaluation, however, a course is an
experimental intervention. Its geometry and topology determine feasible motions; its
width, clearance, visibility, elevation, current, wind, and obstacle semantics shape
the sensing and control problem; and its sampling distribution determines which parts
of that problem a policy encounters. A result obtained on one fixed route, or on an
undisclosed family of procedurally generated routes, cannot establish the same claim
as a result obtained on a versioned distribution and an immutable held-out suite.

This distinction matters at both ends of a learning or control experiment. During
training, generator support, rejection, repair, and curriculum rules determine the
data distribution and can silently remove rare or difficult configurations. During
evaluation, course selection determines measured generalization and sensitivity to
safety-critical geometry. A generator that emits many invalid samples may still
produce an attractive curated set; a controller may achieve high average success
while failing in a narrow clearance or visibility regime; and two simulators may
realize nominally identical centerlines with different collision geometry, frames, or
reset semantics. Without explicit course identities, raw attempt denominators,
vehicle envelopes, and import checks, these cases are difficult to distinguish.

Course generation is consequently more than geometric sampling. A usable artifact
must pass a chain of conditional checks: the generator must terminate; the course must
satisfy declared geometric and topological predicates; it must be feasible for a
stated vehicle, model, and control envelope; and its exported realization must import,
reset, and preserve relevant semantics in the target simulator. None of these stages
is interchangeable. Geometric validity does not prove dynamic traversability, and a
portable file does not prove simulator interoperability. Likewise, controller success
is an outcome conditional on a course, policy, simulator, and seed, not an intrinsic
measure of course quality.

The literature needed to reason about this chain is distributed across several
communities. Surveys of autonomous car and drone racing organize perception,
estimation, planning, control, learning, platforms, and simulators around the racing
system \cite{Betz2022AutonomousVehicles,Hanover2023AutonomousDrone}. Search-based
procedural content generation and virtual-world generation provide transferable ideas
about representations, fitness functions, roads, and world-building pipelines
\cite{Togelius2011SearchBased,Freiknecht2017SurveyVirtualWorlds}. Safety-critical
driving scenario generation contributes methods and criteria for producing test
conditions \cite{Ding2023SafetyCritical}. These perspectives are necessary context,
but their primary units are respectively autonomy systems, game or world content, and
traffic scenarios. As detailed in \cref{sec:adjacent-surveys}, the present survey
instead aligns their transferable concepts around the course artifact and the
evidence required to generate, validate, distribute, import, and evaluate it. This is
a difference in organizing question, not a claim that adjacent surveys contain no
relevant tracks, maps, metrics, or simulators.

The scope also excludes several operations that are easily conflated with course
generation. Racing-line optimization selects a trajectory on an existing course; a
controller maps observations or states to actions; traffic, appearance, weather, and
dynamics randomization alter a broader scenario; and selection, scheduling, replay,
or storage manage existing artifacts. These operations may depend on, transform, or
evaluate courses, but they do not generate a course unless they create or materially
change the traversal-constraining layout. Fixed routes and competition specifications
remain relevant when they expose transferable representations, interfaces, validity
conditions, or benchmark practices, but they are not relabelled as generators.

Against this background, the survey asks three questions. First, how are robot
courses represented and generated across ground, aerial, and maritime systems?
Second, which constraints, validity checks, and evaluation measures make a generated
course usable for learning and control? Third, how should generators, course
distributions, simulator realizations, and policy evaluations be reported so that
results can be reproduced and compared? The unit of technical interest is any layout
that constrains traversal, including closed tracks, open corridors, roads and route
graphs, ordered gates or waypoints, buoy courses, and semantically structured worlds.
The scope is therefore broader than racing geometry alone, but narrower than generic
scenario or world generation.

The paper makes five contributions.
\begin{enumerate}
  \item It defines a cross-domain vocabulary for course objects, traversal relations,
  representations, realizations, generators, validity, training distributions,
  evaluation suites, and controllers. This vocabulary prevents an ordered gate
  sequence, a sampled centerline, a simulator scene, and a policy trajectory from
  being treated as equivalent artifacts.

  \item It synthesizes source-backed representation and generation families together
  with the constraints that change across ground, aerial, and maritime operation.
  Because independent contribution coding did not meet the prespecified reliability
  gate, this synthesis is qualitative; the paper makes no quantitative taxonomy or
  prevalence claims.

  \item It proposes an auditable metric framework that separates generator, course,
  simulator, and policy reporting units. The framework reports feasibility yields
  over all attempts, geometry and topology descriptors, feasible diversity,
  reproducibility and cost, and simulator conformance. It replaces the labels
  \emph{easy}, \emph{medium}, and \emph{hard} with a declared continuous,
  multidimensional descriptor spectrum whose interpretation remains conditional on
  domain and vehicle envelope.

  \item It translates those measures into a prospective benchmark protocol. The
  proposed distribution-estimation and training budget is
  $10{,}000$--$100{,}000$ attempted courses per generator, governed by adaptive
  stopping criteria rather than treated as universally sufficient. A domain and
  configuration enter an immutable evaluation release only after at least 100
  feasible, non-duplicate courses span its declared spectrum and diversity
  constraints; larger suites are required when precision or power demands them.
  Policies are compared on common course identifiers with multiple training and
  rollout seeds and an explicit failure taxonomy.
  The protocol also evaluates each generator as a training distribution through a
  crossed generator-to-suite policy matrix, an evaluation-only real-track column, and
  a separately reported model-based reference-controller calibration.

  \item It specifies a prospective portable \texttt{CourseSpec}, course cards,
  reconstruction and conformance records, and single-environment simulator/RL adapter
  semantics. It then audits the narrower coverage of the current \TrackGen{}
  repository, distinguishes tested, partial, experimental, and planned surfaces, and
  formulates falsifiable hypotheses for future cross-domain evaluation. The portable
  representation, large corpora, benchmark suites, and adapter set are design targets,
  not claims of an existing release or of consensus in the literature.
\end{enumerate}

The review protocol in \cref{sec:review-protocol} separates report-level retention
from work-level technical synthesis and records the limits imposed by incomplete
adjudication and failed contribution-coding reliability. \Cref{sec:scope-definitions}
defines the common objects, and \cref{sec:adjacent-surveys} positions the survey.
\Cref{sec:representations,sec:generation-methods,sec:domain-constraints} synthesize
representations, methods, and domain constraints. \Cref{sec:metrics} defines the
metric contracts; \cref{sec:benchmark-protocol} turns them into a proposed benchmark
and interoperability protocol. \Cref{sec:reference-implementations} records current
implementation coverage and release gaps, \cref{sec:reporting-practices} gives the
reporting standard and audit limits, and \cref{sec:open-problems} states falsifiable
research hypotheses.
